\documentclass[11pt,letterpaper]{article}

% ------ Configuración de Idioma y Codificación ------
\usepackage[utf8]{inputenc}
\usepackage[spanish, es-tabla]{babel} % es-tabla cambia "Cuadro" por "Tabla"

% ------ Paquetes de Color y Estilos Visuales ------
\usepackage[table]{xcolor} % Debe ir arriba para evitar conflictos con tablas
\usepackage{tcolorbox}
\usepackage{epsfig,graphicx}
\usepackage{tikz}

% ------ Configuración de Geometría y Páginas ------
\usepackage[left=2cm,right=2cm,top=1.8cm,bottom=2.3cm]{geometry}
\usepackage{fancyhdr}
\usepackage{lastpage}
\pagestyle{fancy}
\fancyhf{}
\fancyhead[R]{\ifnum\value{page}>1 \textbf{Estadística Bayesiana}\fi}
\lfoot{\includegraphics[height=1.0cm]{ciencias.png}}
\rfoot{\textit{Página \thepage}}

% ------ Paquetes Matemáticos Básicos y Teoremas ------
\usepackage{amsmath}
\usepackage{amssymb}
\usepackage{amsthm}

\newtheorem{theorem}{Theorem}[section]
\newtheorem{corollary}{Corollary}[theorem]
\newtheorem{lemma}[theorem]{Lemma}

% ------ Tablas Avanzadas ------
\usepackage{tabularx}
\usepackage{colortbl}
\usepackage{array, multirow, multicol}

% ------ Formateo de Código (Listings) ------
\usepackage{listings}

% Definición de Colores para el Código
\definecolor{codegreen}{rgb}{0,0.6,0}
\definecolor{codegray}{rgb}{0.5,0.5,0.5}
\definecolor{codepurple}{rgb}{0.58,0,0.82}
\definecolor{backcolour}{rgb}{0.95,0.95,0.92}
\definecolor{maroon}{HTML}{7B1E3A}
\definecolor{AzulUNAM}{RGB}{12, 35, 64}     % Azul marino institucional
\definecolor{OroUNAM}{RGB}{191, 144, 0}     % Oro institucional (el que ya tenías)
\definecolor{backcolour}{rgb}{0.95,0.95,0.92}

% Configuración ÚNICA de mystyle
\lstdefinestyle{mystyle}{
    backgroundcolor=\color{backcolour},
    commentstyle=\color{codegreen},
    keywordstyle=\color{magenta},
    numberstyle=\tiny\color{codegray},
    stringstyle=\color{codepurple},
    basicstyle=\ttfamily\footnotesize,
    breakatwhitespace=false,
    breaklines=true,
    captionpos=b,
    keepspaces=true,
    numbers=left,
    numbersep=5pt,
    showspaces=false,
    showstringspaces=false,
    showtabs=false,
    tabsize=2
}
\lstset{style=mystyle}

% ------ Algoritmos y Pseudocódigo ------
\usepackage{algorithm}
\usepackage{algpseudocode}

% ------ Hipervínculos y Botón de Copiado (Al Final) ------
\usepackage{url}
\usepackage{hyperref}
\usepackage{clipboard} 
% ------ Atajos Matemáticos (Macros) ------
\definecolor{B}{HTML}{FFFFFF}
\definecolor{G}{HTML}{5e5e5e}
\definecolor{R2}{HTML}{d53d40}
\definecolor{A2}{HTML}{034190}
\definecolor{V2}{HTML}{7faa50}

\newcommand{\R}{\mathbb{R}}
\newcommand{\C}{\mathcal{C}}
\newcommand{\N}{\mathbb{N}}
\newcommand{\Z}{\mathbb{Z}}
\newcommand{\Q}{\mathbb{Q}}

% ------ Configuración de Listas Enumeradas ------
\renewcommand{\theenumi}{\Roman{enumi}}
\renewcommand{\labelenumi}{\theenumi.}

\begin{document}

%---Evironment--%
	\makeatletter
        \renewenvironment{proof}[1][\proofname]{\par
            \pushQED{\qed}%
            \normalfont \topsep6\p@\@plus6\p@\relax
            \trivlist
            \item\relax
            {\itshape
            #1\@addpunct{.}}\par\vspace{\baselineskip}\ignorespaces
            }{%
            \popQED\endtrivlist\@endpefalse
            }
        \newenvironment{Solucion}[1][Soluci\'on]{\par
            \pushQED{\qed}%
            \normalfont \topsep6\p@\@plus6\p@\relax
            \trivlist
            \item\relax
            {\itshape
            #1\@addpunct{.}}\par\vspace{\baselineskip}\ignorespaces
            }{%
            \popQED\endtrivlist\@endpefalse
            }
            
    \makeatother

%------ Encabezado -------- %

\begin{center}
    \begin{minipage}{3cm}
    	\begin{center}
    		\includegraphics[height=3.5cm]{unam.jpg}
    	\end{center}
    \end{minipage}\hfill
    \begin{minipage}{10cm}
    	\begin{center}
    	\textbf{\large Universidad Nacional Autónoma de México}\\[0.1cm]
        \textbf{Facultad de Ciencias}\\[0.1cm]
        \textbf{\'Estadistica Bayesiana}\\[0.1cm]
        Tarea grande 1 \\[0.1cm]
        El\'ias L\'opez Rivera\\[0.1cm]
        Hector Gonzalez Baltierra
        \\[0.1cm]
        %\texttt{ elias.lopezr,@ciencias.unam.mx }\\[0.1cm]
        Fecha:\,\,24/09/2026
    	\end{center}
    \end{minipage}\hfill
    \begin{minipage}{3cm}
    	\begin{center}
    		\includegraphics[height=3.5cm]{ciencias.png}
    	\end{center}
    \end{minipage}
\end{center}
\rule{17cm}{0.1mm}
%------ Fin de encabezado -------- %
\begin{tcolorbox}[
    title = \textcolor{white}{\textbf{Problema 1}},
    colbacktitle=AzulUNAM!95!black,         % Fondo del título en azul marino profundo
    coltitle=white,                         % Texto del título en blanco
    colback=OroUNAM!7!white,                % Fondo del cuerpo con un ligero toque de oro
    colframe=OroUNAM!90!black,              % Borde del recuadro en oro institucional
    fontupper=\small
]
\textit{Un médico cree que su paciente tiene una y solo una de las enfermedades \(E_1, E_2, E_3, E_4\) y que sus probabilidades respectivas son
\[
p(E_1) = 0.2, \; p(E_2) = 0.5, \; p(E_3) = 0.2 \; \text{ y } \; p(E_4) = 0.1.
\]
El médico dispone de tres tratamientos alternativos, \(\{t_1, t_2, t_3\}\), cuya efectividad viene reflejada en la siguiente tabla:
\begin{center}
%\renewcommand{\arraystretch}{1.2}
\begin{tabular}{c|cccc}
 & \(E_1\) & \(E_2\) & \(E_3\) & \(E_4\) \\ \hline
\(t_1\) & 1 & 0 & 1 & 1 \\
\(t_2\) & 0 & 1 & 0 & 1 \\
\(t_3\) & 0 & 1 & 1 & 0 \\
\end{tabular}
\end{center}
en donde un uno en la fila \(i\) y columna \(j\) significa que el tratamiento \(t_i\) es efectivo contra la enfermedad \(E_j\) y un 0 significa que no lo es. Igualando utilidad a efectividad, determinen el tratamiento óptimo. Ahora consideren que si un tratamiento no es efectivo, puede procederse a otro y determinen la sucesión óptima de tratamientos, suponiendo que los costos respectivos son \(c_1 = 0.1, c_2 = 0.2, c_3 = 0.3\) en unidades de utilidad.
}
\end{tcolorbox}

\begin{proof}

    Sea E el conjunto de estados:

    \[
    E=\{E_1,E_2,E_3,E_4\}.
    \]

    Sea T el conjunto de tratamientos:

    \[
    T=\{t_1,t_2,t_3\}.
    \]

    El conjunto de resultados:

    \[
    R=\{\text{efectivo},\text{no efectivo}\}.
    \]

    Por hipótesis la distribución de probabilidad sobre los estados es:

    \[
    \begin{aligned}
    \pi(E_1)&=0.2, & \pi(E_2)&=0.5,\\ \pi(E_3)&=0.2, & \pi(E_4)&=0.1.
    \end{aligned}
    \]

    Definimos la función de efectividad como

    \[
    e:T\times E\xrightarrow[]{}\{0,1\},
    \]

    $e(t_i,E_j)=1$ si el tratamiento $t_i$ es efectivo con $E_j$, y $e(t_i,E_j)=0$ si no. Usaremos esta matriz
    
    \[
    \begin{pmatrix}
    1&0&1&1\\
    0&1&0&1\\
    0&1&1&0
    \end{pmatrix},
    \]

    Definimos la función de costos

    \[
    c:T\xrightarrow[]{}\mathbb{R},
    \]

    por

    \[
    c(t_1)=0.1,\ \ \ c(t_2)=0.2,\ \ \ c(t_3)=0.3.
    \]

    Consideraremos el caso en el que no hay costos:

    \[
    \mathbb{E}_\pi[e(t_i,E)]=\sum_{j=1}^{4}\pi(E_j)e(t_i,E_j).
    \]

    Así tenemos

    \[
    \begin{aligned}
    \mathbb{E}_\pi[e(t_1,E)] &=1(0.2)+0(0.5)+1(0.2)+1(0.1)\\ &=0.2+0.2+0.1\\ &=0.5.
    \end{aligned}
    \]

    \[
    \begin{aligned}
    \mathbb{E}_\pi[e(t_2,E)] &=0(0.2)+1(0.5)+0(0.2)+1(0.1)\\ &=0.5+0.1\\ &=0.6.
    \end{aligned}
    \]
    
    \[
    \begin{aligned}
    \mathbb{E}_\pi[e(t_3,E)] &=0(0.2)+1(0.5)+1(0.2)+0(0.1)\\ &=0.5+0.2\\ &=0.7.
    \end{aligned}
    \]
    
    Así,

    \[
    \operatorname*{arg\,max}_{t_i\in T}\mathbb{E}_\pi[e(t_i,E)]=\{t_3\}.
    \]

    Si no consideramos costos el tratamiento efectivo es $t_3$.\\



   Si consideramos costos y además el poder aplicar un tratamiento si no funciona el anterior entonces es necesario ampliar nuestro modelo \\

    Para \(k\in\{1,2,3\}\), definimos

    \[
    V_k=\left\{(t_{i_1},\ldots,t_{i_k})\in T^k:i_r\neq i_s\text{ si }r\neq s\right\}.
    \]

    Definimos

    \[
    V=V_1\cup V_2\cup V_3.
    \]

    Definimos la función de utilidad extendida
    
    \[
    u:V\times E\xrightarrow[]{}\mathbb{R}.
    \]

    Para una sucesión \(\sigma=(t_{i_1},\ldots,t_{i_k})\in V_k\) y un estado \(E_l\), su utilidad es

    Si k=1

    \[
    u(\sigma,E_l)=
    \begin{cases}
    1-c(t_{i_1}), & e(t_{i_1},E_l)=1,\\
    -c(t_{i_1}), & 1-e(t_{i_1},E_l)=1.
    \end{cases}
    \]

    Si k=2

    \[
    u(\sigma,E_l)=
    \begin{cases}
    1-c(t_{i_1}),
    & \text{si }e(t_{i_1},E_l)=1,\\
    1-c(t_{i_1})-c(t_{i_2}),
    & \text{si }(1-e(t_{i_1},E_l))=1,\ e(t_{i_2},E_l)=1,\\
    -c(t_{i_1})-c(t_{i_2}),
    & \text{si }(1-e(t_{i_1},E_l))(1-e(t_{i_2},E_l))=1.
    \end{cases}
    \]

    Si k=3

    \[
    u(\sigma,E_l)=
    \begin{cases}
    1-c(t_{i_1}),
    & \text{si }e(t_{i_1},E_l)=1,\\
    1-c(t_{i_1})-c(t_{i_2}),
    & \text{si }(1-e(t_{i_1},E_l))=1,\ e(t_{i_2},E_l)=1,\\
    1-c(t_{i_1})-c(t_{i_2})-c(t_{i_3}),
    & \text{si }(1-e(t_{i_1},E_l))(1-e(t_{i_2},E_l))=1,\ e(t_{i_3},E_l)=1,\\
    -c(t_{i_1})-c(t_{i_2})-c(t_{i_3}),
    & \text{si }(1-e(t_{i_1},E_l))(1-e(t_{i_2},E_l))(1-e(t_{i_3},E_l))=1.
    \end{cases}
    \]

    \[
    \mathbb{E}_\pi[u(\sigma,E)]=\sum_{l=1}^{4}\pi(E_l)u(\sigma,E_l).
    \]

    Así

    \[
    \begin{aligned}
    \mathbb{E}_\pi[u((t_1),E)] &=0.2(1-0.1)+0.5(0-0.1)\\ &+0.2(1-0.1)+0.1(1-0.1)\\ &=0.4,
    \end{aligned}
    \]
    
    \[
    \begin{aligned}
    \mathbb{E}_\pi[u((t_2),E)] &=0.2(0-0.2)+0.5(1-0.2)\\ &+0.2(0-0.2)+0.1(1-0.2)\\ &={0.4},
    \end{aligned}
    \]

    \[
    \begin{aligned}
    \mathbb{E}_\pi[u((t_3),E)] &=0.2(0-0.3)+0.5(1-0.3)\\ &+0.2(1-0.3)+0.1(0-0.3)\\ &={0.4}.
    \end{aligned}
    \]
    
    \[
    \begin{aligned}
    \mathbb{E}_\pi[u((t_1,t_2),E)] &=0.2(1-0.1)+0.5(1-0.1-0.2)\\ &+0.2(1-0.1)+0.1(1-0.1)\\ &={0.8},
    \end{aligned}
    \]
    
    \[
    \begin{aligned}
    \mathbb{E}_\pi[u((t_1,t_3),E)] &=0.2(1-0.1)+0.5(1-0.1-0.3)\\ &+0.2(1-0.1)+0.1(1-0.1)\\ &={0.75},
    \end{aligned}
    \]
    
    \[
    \begin{aligned}
    \mathbb{E}_\pi[u((t_2,t_1),E)]  &=0.2(1-0.2-0.1)+0.5(1-0.2)\\ &+0.2(1-0.2-0.1)+0.1(1-0.2)\\ &={0.76},
    \end{aligned}
    \]
    
    \[
    \begin{aligned}
    \mathbb{E}_\pi[u((t_2,t_3),E)] &=0.2(0-0.2-0.3)+0.5(1-0.2)\\ &+0.2(1-0.2-0.3)+0.1(1-0.2)\\ &={0.48},
    \end{aligned}
    \]
    
    \[
    \begin{aligned}
    \mathbb{E}_\pi[u((t_3,t_1),E)] &=0.2(1-0.3-0.1)+0.5(1-0.3)\\ &+0.2(1-0.3)+0.1(1-0.3-0.1)\\ &={0.67},
    \end{aligned}
    \]
    
    \[
    \begin{aligned}
    \mathbb{E}_\pi[u((t_3,t_2),E)] &=0.2(0-0.3-0.2)+0.5(1-0.3)\\ &+0.2(1-0.3)+0.1(1-0.3-0.2)\\ &={0.44}.
    \end{aligned}
    \]
    
    \[
    \begin{aligned}
    \mathbb{E}_\pi[u((t_1,t_2,t_3),E)] &=0.2(1-0.1)+0.5(1-0.1-0.2)\\ &+0.2(1-0.1)+0.1(1-0.1)\\ &={0.8},
    \end{aligned}
    \]
    
    \[
    \begin{aligned}
    \mathbb{E}_\pi[u((t_1,t_3,t_2),E)] &=0.2(1-0.1)+0.5(1-0.1-0.3)\\ &+0.2(1-0.1)+0.1(1-0.1)\\ &={0.75},
    \end{aligned}
    \]
    
    \[
    \begin{aligned}
    \mathbb{E}_\pi[u((t_2,t_1,t_3),E)] &=0.2(1-0.2-0.1)+0.5(1-0.2)\\ &+0.2(1-0.2-0.1)+0.1(1-0.2)\\ &={0.76},
    \end{aligned}
    \]
    
    \[
    \begin{aligned}
    \mathbb{E}_\pi[u((t_2,t_3,t_1),E)] &=0.2(1-0.2-0.3-0.1)+0.5(1-0.2)\\ &+0.2(1-0.2-0.3)+0.1(1-0.2)\\ &={0.66},
    \end{aligned}
    \]
    
    \[
    \begin{aligned}
    \mathbb{E}_\pi[u((t_3,t_1,t_2),E)] &=0.2(1-0.3-0.1)+0.5(1-0.3)\\ &+0.2(1-0.3)+0.1(1-0.3-0.1)\\ &={0.67},
    \end{aligned}
    \]
    
    \[
    \begin{aligned}
    \mathbb{E}_\pi[u((t_3,t_2,t_1),E)] &=0.2(1-0.3-0.2-0.1)+0.5(1-0.3)\\ &+0.2(1-0.3)+0.1(1-0.3-0.2)\\ &={0.62}.
    \end{aligned}
    \]

    \[
    \max_{\sigma\in V}\mathbb{E}_\pi[u(\sigma,E)]=0.8.
    \]

    \[
    \operatorname*{arg\,max}_{\sigma\in V}\mathbb{E}_\pi[u(\sigma,E)]=\{(t_1,t_2)(t_1,t_2,t_3)\}.
    \]


\end{proof}\,

\begin{tcolorbox}[
    title = \textcolor{white}{\textbf{Problema 2}},
    colbacktitle=AzulUNAM!95!black,         % Fondo del título en azul marino profundo
    coltitle=white,                         % Texto del título en blanco
    colback=OroUNAM!7!white,                % Fondo del cuerpo con un ligero toque de oro
    colframe=OroUNAM!90!black,              % Borde del recuadro en oro institucional
    fontupper=\small
]
\textit{Supongan que se sabe que $\theta$ es mayor que cero y que la incial es
uniforme en ese rango. Una muestra aleatoria de tamaño $16$ de una
$N (x|\theta,1)$, arrojó una media de $0.40$. Obtengan un intervalo de proba-
bilidad $0.87$ para $\theta$.}
\end{tcolorbox}

\begin{tcolorbox}[
    title = \textcolor{white}{\textbf{Problema 3}},
    colbacktitle=AzulUNAM!95!black,         % Fondo del título en azul marino profundo
    coltitle=white,                         % Texto del título en blanco
    colback=OroUNAM!7!white,                % Fondo del cuerpo con un ligero toque de oro
    colframe=OroUNAM!90!black,              % Borde del recuadro en oro institucional
    fontupper=\small
]
\textit{Dos muestras aleatorias fueron tomadas de distribuciones normales
con medias distintas, pero con la misma precisi\'on. Encuentre la dis-
tribuci\'on final de la precisi\'on, suponiendo que la información obtenida
fue $\{9.1, 9.3, 9.9, 8.8, 9.4\}$ para el primer caso y $\{10.6, 10.4, 9.3, 10.1,
10.1, 10.0, 10.7, 9.6\}$ para el segundo. Supongan que $\pi(\mu_1,\mu_2,h) \propto h^{-1}$
.}
\end{tcolorbox}
\begin{tcolorbox}[
    title = \textcolor{white}{\textbf{Problema 4}},
    colbacktitle=AzulUNAM!95!black,         % Fondo del título en azul marino profundo
    coltitle=white,                         % Texto del título en blanco
    colback=OroUNAM!7!white,                % Fondo del cuerpo con un ligero toque de oro
    colframe=OroUNAM!90!black,              % Borde del recuadro en oro institucional
    fontupper=\small
]
\textit{\begin{itemize}
    \item Desde el punto de vista frecuentista, un intervalo de confianza, digamos $90\%$ para el parametro $\theta$ de una Bernoulli, se construye de la siguiente manera:
    \begin{enumerate}
        \item Se considera el estimador m\'aximo veros\'imil de $\theta$, la media aritm\'etica.
        \item Se supone que (el tamaño de la muestra deberia ser grande)
        \begin{equation*}
            \widehat{\theta}\sim N\left(\widehat{\theta}|\theta,\frac{\theta(1-\theta)}{n}\right) 
        \end{equation*}
        \item Como $\widehat{\theta}\to \theta $, entonces $\theta(1-\theta)\approx\widehat{\theta}(\widehat{\theta}-1)  $
        y entonces proponen como un intervalo de confianza a:\,\\
        \begin{equation*}
            \left(\widehat{\theta}-1.64\sqrt{\frac{\widehat{\theta}(\widehat{\theta}-1)}{n}},\widehat{\theta}+1.64\sqrt{\frac{\widehat{\theta}(\widehat{\theta}-1)}{n}}\right)
        \end{equation*}
    \end{enumerate}
        Esto significa que si construimos un gran n\'umero de intervalos de la misma forma, el $90\%$ de ellos contendran al verdadero valor del par\'ametro.¿Ser\'a cierto?. Hagan un pequeño 
        ejercicio de simulaci\'on para verificarlo:
        \begin{enumerate}
            \item Supongan un valor fijo de $\theta$
            \item Generen $1000$ muestras  de tamaño $n=20$ de una Bernoulli con este valor de $\theta$
            \item Construyan los $10,000$ intervalos de confianza y cuenten cu\'antas veces el intervalo contiene el valor $0.35$
        \end{enumerate}
        Este porcentaje se llamam cobertura frecuentista, y si las cosas estan bien, la ley de los grande n\'umeros les dice que, en este caso, ese valor debe ser cercano a $0.90$. Comenten sus resultados
        \item Si suponemos $\pi(\theta)\propto \theta^{\frac{1}{2}}(1- \theta)^{-\frac{1}{2}}$, la distribuci\'on final de $\theta$ ser\'a una beta. Use las muestras obtenidas en la primera parte para calsular $1000$ intervalos de longitud m\'inima basados en
        la distribuci\'on final correspondiente.el porcentaje de veces que $\theta$ est\'a en los intervalos ser\'a la correspondiente cobertura frecuentista
        \item ¿La cobertura depende del tamaño de muestra? Una forma de verlo, es repetir lo anterior para diferentes tamaños d emuestra tales como $\{20,30,40,50,60,70,80,90,100\}$
        \item ¿La cobertura dependera del valor de $\theta$?.Repitan todo para $\theta\in \{0.05,0.20,0.5\}$
    \end{itemize}
}
\end{tcolorbox}
\begin{tcolorbox}[
    title = \textcolor{white}{\textbf{Problema 5}},
    colbacktitle=AzulUNAM!95!black,         % Fondo del título en azul marino profundo
    coltitle=white,                         % Texto del título en blanco
    colback=OroUNAM!7!white,                % Fondo del cuerpo con un ligero toque de oro
    colframe=OroUNAM!90!black,              % Borde del recuadro en oro institucional
    fontupper=\small
]
\textit{Suponga que $\textbf{Y}=\{y_1, . . . , y_m\}$ es una muestra aleatoria de una
distribuci\'on uniforme en el intervalo $(0,\infty)$. Sea $\gamma= log\theta$ y suponga
que la distribuci\'on inicial de $\gamma$ es normal com media $\mu$ y varianza $\sigma^2$.Encuentre la distribuci\'on final de $\gamma$ dado $\textbf{Y}$, y muestre que es una
distribuci\'on normal truncada.}
\end{tcolorbox}
\begin{proof}
    Tenemos que $p(\gamma)\propto e^{-\frac{(\gamma-\mu)^2}{2\sigma^2}}$
    ademas $p(\gamma|\textbf{Y})\propto p(\textbf{Y}|\gamma)p(\gamma)$, obtendremos la verosimilitud $p(\textbf{Y}|\gamma)$:\,\\
    \begin{equation*}
        p(\textbf{Y}|\gamma)=\prod_{i=1}^{m}P(y_i|\gamma)
    \end{equation*}\,\\
    Es claro que $P(y_i|\theta)=\frac{1}{\theta}\mathbb{I}(0<y<\theta)$, por tanto:\,\\
    \begin{equation*}
        p(\textbf{Y}|\theta)=\prod_{i=1}^m p(y_i|\theta)=\prod_{i=1}^m \frac{1}{\theta}\mathbb{I}(0<y_i<\theta)=\prod_{i=1}^m\theta^{-m}\,\mathbb{I}(\theta>y_{(m)})
    \end{equation*}\,\\
    \,\\
    Donde $y_{(m)}=max\{y_1,y_2,\cdots,y_m\}$,despejando $\gamma=log(\theta)\implies \theta=e^{\gamma}$:\,\\
    \begin{equation*}
        p(\textbf{Y}|\gamma)=\prod_{i=1}^m e^{-m\gamma}\,\mathbb{I}(e^{\gamma}>y_{(m)})=\prod_{i=1}^m e^{-m\gamma}\,\mathbb{I}(\gamma>log(y_{(m)}))
    \end{equation*}
    Por tanto:\,\\
    \begin{equation*}
        p(\gamma|\textbf{Y})\propto p(\textbf{Y}|\gamma)p(\gamma)=e^{-\frac{(\gamma-\mu)^2}{2\sigma^2}-m\gamma}\,\,\mathbb{I}(\gamma>log(\theta))
    \end{equation*}\,\\
    Expandiendo  y completando el cudraro:\,\\
    \begin{align*}
        -\frac{1}{2\sigma^2}(\gamma^2-2\gamma\mu+\mu^2)-m\gamma&=-\frac{1}{2\sigma^2}(\gamma^2-2(\mu-m\sigma^2)+\mu^2)\,\\
        \,\\
        &=-\frac{1}{2\sigma^2}(\gamma^2-2(\mu-m\sigma^2)\gamma+(\mu-m\sigma^2)^2+\mu^2-(\mu-m\sigma^2)^2)
        \,\\
        \,\\
        &=-\frac{1}{2\sigma^2}((\gamma-(\mu-m\sigma^2))^2+\mu^2-(\mu-m\sigma^2)^2)
    \end{align*}\,\\
    Por tanto $e^{-\frac{(\gamma-\mu)^2}{2\sigma^2}-m\gamma}=e^{\mu^2-(\mu-m\sigma^2)^2}e^{-\frac{(\gamma-(\mu-m\sigma^2))^2}{2\sigma^2}}$, es decir:\,\\
    \begin{equation*}
         p(\gamma|\textbf{Y})\propto p(\textbf{Y}|\gamma)p(\gamma)\propto e^{-\frac{(\gamma-(\mu-m\sigma^2))^2}{2\sigma^2}}\,\,\mathbb{I}(\gamma>log(\theta))
    \end{equation*}\,\\
    \,\\
    Que es el kernel de una Normal truncada.
\end{proof}\,\\


\begin{tcolorbox}[
    title = \textcolor{white}{\textbf{Problema 6}},
    colbacktitle=AzulUNAM!95!black,         % Fondo del título en azul marino profundo
    coltitle=white,                         % Texto del título en blanco
    colback=OroUNAM!7!white,                % Fondo del cuerpo con un ligero toque de oro
    colframe=OroUNAM!90!black,              % Borde del recuadro en oro institucional
    fontupper=\small
]
\textit{Sea $X$ una caracter\'istica (univariada) que puede medirse en una po
blaci\'on finita. Por ejemplo, $X$ es la edad de un conjunto de personas de
una cierta poblaci\'on, o la preferencia partidista o el ingreso promedio,
etc. Supongamos que la poblaci\'on consta de $m$ individuos, as\'i que a lo
m\'as podemos observar ${X1, . . . , Xm}$.
\begin{itemize}
    \item Supongamos que la distribuci\'on inicial $\{X_1,\cdots,X_m\}$ podemos escribirla como:\,\\
    \begin{equation*}
        P(X_1,\cdots,X_m)=\int \prod_{i=1}^m p(X_i|\theta)p(\theta)d\theta
    \end{equation*}
    para algunas $p(X|\theta)$ y $p(\theta)$. Si observamos una muestra de $\{X_1,\cdots,X_m\}$, digamos $\{x_1,\cdots,x_m\}$ con $n<m$
    encuentre la distribuci\'on final:\,\\
    \begin{equation*}
        p(X_1,\cdots,X_m|x_1,\cdots,x_m)
    \end{equation*}
    en t\'erminos de la densidad condicional de $X$ dado $\theta$, la inicial de $\theta$ y la verosimilitud observada
    \item Si en el inciso anterior, suponemos que la densidad condicional de $X$ dado $\theta$ es una normal con media $\mu$ y presici\'on $1$, y la distribuci\'on inicial para $\theta$ es una normal con par\'ametros
    $(\mu,\tau)$ ($\tau$ presici\'on), encuentre:\,\\
    \begin{equation*}
        P(X_1|x_1,\cdots,x_n)
    \end{equation*}
    cuando:
    i).\,$X_1$ est\'a en la muestra
    ii).\,$X_1$ no est\'a en la muestra
\end{itemize}}
\end{tcolorbox}
\begin{tcolorbox}[
    title = \textcolor{white}{\textbf{Problema 7}},
    colbacktitle=AzulUNAM!95!black,         % Fondo del título en azul marino profundo
    coltitle=white,                         % Texto del título en blanco
    colback=OroUNAM!7!white,                % Fondo del cuerpo con un ligero toque de oro
    colframe=OroUNAM!90!black,              % Borde del recuadro en oro institucional
    fontupper=\small
]
\textit{Sup\'ongase que $X_i\sim Po(x_i|\theta_i)$, $i\in J_m$ y que estas variables aleatorias son condicionalmente independientes dados los valores de los par\'ame-
tros. Sup\'ongase adem\'as que $\theta_1, . . . , \theta_m$ se distribuyen de acuerdo a una
exponencial con media $\phi^{-1}$ y que son condicionalmete independientes
dado el valor de $\phi$. Finalmente se considera una distribuci\'on inicial
para $\phi$ tak que $p(\phi) \propto1$.
\begin{itemize}
    \item Sea $\varphi=(1+\phi)^{-1}$. Demuestra que la distribuci\'on final de $\varphi$ es proporcional a:\,\\
    \begin{equation*}
        \phi^{m\bar{x}-2}(1+\phi)^n
    \end{equation*} 
    \item Usando este resultado, encuentre el estimador de $\theta_i$, $i\in J_m$ basado en una funci\'on de p\'erdida cuadr\'atica
    \item ¿Qu\'e sucede si $m\bar{x}\leq 1$ 
\end{itemize}}
\end{tcolorbox}
\begin{proof}
    Notemos que como $p(\phi|X)\varpropto p(X|\phi)p(\phi)\varpropto p(X|\phi) $ primero encontraremos $p(X|\phi)$:\,\\
    \begin{equation*}
        p(X|\phi)=\int_{0}^{\infty}\cdots\int_{0}^{\infty} p(X|\theta)p(\theta|\phi)d\theta_1\cdots d\theta_m
    \end{equation*}
    Donde $\theta=(\theta_1,\cdots,\theta_m)$ es el vector correpodiente de par\'ametros, usando la independencia condicional de los $\theta$'s con $\phi$ y que la muestra $X$ tambi\'en es condicionalmente independiente tenemos que:\,\\
    \begin{equation*}
        p(X|\phi)=\prod_{i=1}^m\int_{0}^{\infty}p(x_i|\theta_i)p(\theta_1|\phi)d\theta_i
    \end{equation*}\,\\
    \,\\
    Finalmente como $X_i\sim po(x_i|\theta_i)$ y que $\theta_i$ de sistribuye con una exponencial de media $\phi^{-1}$:\,\\
    \,\\
     \begin{align*}
        p(X|\phi)&=\prod_{i=1}^m\int_{0}^{\infty}\frac{\theta_i^{x_i}e^{-\theta_i}}{x_i!}\phi e^{-\phi\theta_i} d\theta_i\,\\
        \,\\
        &=\prod_{i=1}^m\frac{\phi}{x_i!}\prod_{i=1}^m\int_{0}^{\infty}\theta_i^{(x_i+1)-1} e^{-(1+\phi)\theta_i}d\theta_i
    \end{align*}\,\\
    \,\\
    Reconocemos inmediatamente el kernel de una $Ga(x_i+1,(1+\phi)^{-1})$, por tanto:\,\\
    \,\\
    \begin{equation*}
        p(X|\phi)=\prod_{i=1}^m\frac{\phi}{x_i!}\frac{\Gamma(x_i+1)}{(1+\phi)^{x_i+1}}=\frac{\phi\,\,x_i!}{x_i!(1+\phi)^{x_i+1}}=\frac{\phi^n}{(1+\phi)^{\sum_{i=1}^m(x_i+1)}}
    \end{equation*}\,\\
    \,\\
    Renombramos $\sum_{i=1}^mx_i=m\bar{x}$ por tanto:\,\\
    \begin{equation*}
        p(\phi|X)\propto p(X|\phi)=\frac{\phi^n}{(1+\phi)^{m\bar{x}+m}}
    \end{equation*}\,\\
    \,\\
    Finalmente hacemos el cambio de varia $\varphi=(1+\phi)^{-1}$ por tanto:\,\\
    \,\\
    \begin{equation*}
       p(\varphi|X)=p(\phi|X)\left|\frac{d\phi}{d\varphi}\right|\propto \frac{\left(\frac{1-\varphi}{\varphi}\right)^m}{\left(\frac{1}{\phi}\right)^{m\bar{x}+m}}\frac{1}{\varphi^2}=(1-\varphi)^m\frac{\varphi^{m\bar{x}+m}}{\varphi^m} \varphi^{-2}=(1-\varphi)^{m}\varphi^{m\bar{x}-2}
    \end{equation*}\,\\
    \,\\
    Ahora como ya sabemos el estimador para $\theta_i$ que minimiza la perdida cuadratica es $widehat{\theta_i}=E[\theta_i|X]$, usando la ley de esperanza total:\,\\
    \begin{equation*}
        E[\theta_i|X]=E[E[\theta_i|X,\varphi]|x]
    \end{equation*}
    Encontramos la distribucion de $p(\theta_i|x,\varphi)$ notemos que como como la muestra $X$ son variables aleatorias independientes $\theta_i$ solo depende de $x_i$ por tanto:\,\\
    \,\\
    \begin{equation*}
        p(\theta|x_i,\varphi)\propto p(x_i|\theta_i)p(\theta_i|\varphi)\propto \theta_i^{x_i}e^{-\theta_i}e^{-\varphi\theta_i}=\theta_i^{x_i}e^{-(1+\varphi)\theta_i}
    \end{equation*}\,\\
    \,\\
    Por tanto $\theta_i|x_i,\varphi\sim Ga(x_i+1,\varphi^{-1})$, es decir:\,\\
    \,\\
    \begin{equation*}
        E[\theta_i|x,\varphi]=(x_i+1)\varphi\implies \widehat{\theta_i}=E[E[\theta_i|X,\varphi]|x\,]=(x_i+1)E[\varphi|x] 
    \end{equation*}\,\\
    \,\\
    Ya demostramos que $\varphi|X\sim Beta(m\bar{x}-1,m+1)$ es decir:\,\\
    \,\\
    \begin{equation*}
        \widehat{\theta_i}=(x_i+1)E[\phi|X\,]=(x_i+1)\frac{m\bar{x}-1}{(m\bar{x}-1)+(m+1)}=(x_i+1)\frac{m\bar{x}-1}{m(\bar{x}+1)} 
    \end{equation*}

\end{proof}
\end{document}
